% =============================================================================
%  \section{Methodology}  --  PAPER VERSION, IEEEtran conference, 6 pages
%  Drafted 2026-09-24, from the block parked at the end of
%  sec_system_model_waveform.tex, checked against cgan/{train_gan,
%  train_shaped,attacks,detectors,channel}.py and README Sec. 2.10, 3.1,
%  3.3e, 3.3f. Pairs with sec_system_model_waveform.tex and uses its labels
%  sec:system, sec:attacker, sec:threat, eq:power, eq:problem, eq:gan.
%
%  WHAT THIS DESCRIBES
%  -------------------
%  ONLY the method: how the generator G of Sec. attacker is trained against
%  the defender -- initialization by Zhou's imitation training (one cited
%  sentence; the architecture itself is parked for Setup since the 2026-09-26
%  reframe), then training through the frozen detector. What the comparisons isolate, and the
%  hard-decision reporting rule, moved to Experiment Setup 2026-09-26 (user:
%  experimental design, not method; parked below). The 2026-09-24 split (user):
%    * System Model = WHAT the jammer can transmit (abstract G, Eq. eq:gan,
%      tier T1); Methodology = HOW G is chosen. A sentence that would change
%      if G were built or trained differently belongs here;
%    * the baseline attacks (barrage, pulsed, shaped control, genie) and the
%      ladder table are in Experiment Setup -- so the shaped control's CMA-ES
%      recipe is parked below for Experiment Setup, next to its definition.
%  NOT values, NOT the evaluation protocol, NOT results (all parked below).
%  No longer \ref's sec:setup in the text, but the parked blocks do: add \label{sec:setup} to Overleaf's
%  \section{Experiment Setup} when pasting.
%
%  VALUES  (for Experiment Setup; none belong in this section's text -- since
%  the reframe not even Zhou's stated dimensions z 400, label 16, M 1024)
%  ------------------------------------------------------------------
%   plain G (run001_G, train_cgan.py): channels 64/128/256, kernel 5, dropout
%        0.3, pool-to-16; nsgan + GP 10, n_critic 1, Adam 2e-4 (0.5, 0.999),
%        all loss weights 1, 10 000 iterations, batch 128.
%   detector-aware G (train_gan.py --init random): random init (run003 --
%        the paper numbers since 2026-09-27; run002 = the same from the
%        imitation start, run001 = imitation start at 400 steps, both in the
%        appendix); Adam 2e-4 (0.5, 0.999); 4000 steps;
%        128 frames/step (32 for the CNN); JSR band
%        (-32, 0) dB for power/kurtosis, (-48, -16) dB for the CNN;
%        beta in {1, 10, 100, 1000}, CNN also {1e4, 1e5, 1e6}; thresholds at
%        alpha = 0.05; soft-P_det width = clean std over 4096 frames.
%        20 generators in all: 19 detector-aware + the beta = 0 one (task 0).
%   parameter count: G = 8,597,136 (run001_G.pt, counted 2026-09-26; 8.39 M
%        of it is the 4096->2048 dense layer) vs the control's 48. The README
%        said "1M" until 2026-09-26 -- wrong, corrected there too.
%   REPRODUCIBILITY: run001_G's recipe (nsgan, MSE feature matching, log1p
%        STFT, headroom 1.0) is NOT what train_cgan.py's current defaults
%        produce (WGAN-GP, L1, log-clamp, headroom 0.95). The text below is
%        generic enough to cover both; Setup must quote run001's.
%
%  NOTATION  (kept clear of the System Model's symbols)
%  --------
%   psi(r) = a detector's scalar statistic (T is taken by the tiers T0-T2,
%   sigma by the noise std); tau_alpha = its threshold; s_psi = its clean std.
%   beta keeps its System Model meaning, but on the log-BER scale of
%   Eq. (loss) its numerical range differs -- the text says so.
%
%  CITE KEYS: zhou2025cgan (Overleaf key 11184988 -- rename one side),
%  bengio2013estimating (NEW, in paper_drafts/refs_new.bib).
%  Parked text: hansen2016cma (NEW, same file).
%
%  LENGTH: ~500 words of text, about 0.8 column (the System Model is ~1,050,
%  ~1.6 columns). If over budget, cut in this order: the JSR-band sentence
%  (can move to Experiment Setup with its values); the one-generator-per-
%  target sentence; only then touch the surrogates.
%  Register follows Zhang & Krunz (MILCOM'23): the method names its knobs,
%  the evaluation section gives their values.
%
%  REFRAMED 2026-09-26 on the supervisor's input (README Sec. B.2, 2.9): "the
%  main contrib now seems the application of some neural architecture ...
%  novelty in how defenders and attackers interact, coordinate, and
%  syncronize". The section now leads with HOW THE ATTACKER INTERACTS WITH THE
%  DEFENDER -- training through the deployed detector -- and the architecture
%  is one cited sentence; its description is parked for Experiment Setup
%  below. Team synchronization (D4a, README Sec. 3.3i) goes in only if time
%  allows (user); until then it is parked as Future Work below.
%
%  METHOD CHANGED 2026-09-27 (user: "change ... methodology to new and better
%  method, we want best results"): the generator is trained FROM RANDOM
%  WEIGHTS, 4000 steps -- no imitation stage. Zhou is cited for the
%  architecture only. Why: at matched BER the imitation start is detected far
%  more by the CNN (0.48 vs 0.15, 4 seeds, 30 dB) because it biases G toward
%  the victim's modulation (README Sec. 3.3f). The imitation stage and the
%  400-step budget are in the Experiment History appendix
%  (paper_drafts/sec_exphist.tex), promotable to the main text if space
%  allows (user). Paper numbers = run003 (README Sec. 3.3f).
% =============================================================================

\section{Methodology}
\label{sec:method}
The jammer learns to evade the defender by training against the defender's own detection pipeline.
At each training step, its waveform passes through the link of Eq.~\eqref{eq:gan} and is scored twice, by the damage it does to the victim's decisions and by the response it draws from the frozen detector; both are back-propagated into the generator.
No other training signal is used.

\subsection{Generator}
\label{sec:method:plain}
% Was "Initialization by Imitation" until 2026-09-27; label kept. Architecture
% details parked for Experiment Setup below. cgan/models.py; train_gan.py
% --init random.
The generator is a generic convolutional network with a dense output layer that maps $\mathbf{z}$ to one segment of $M$ complex samples; we take its layout from the conditional generative adversarial network (GAN) of~\cite{zhou2025cgan}, but not its training, and nothing in the method depends on it.
% User 2026-09-27. Zhou cites no source for the generator: noise (400) + label
% embedding (16) -> 3 Conv1d blocks over the noise vector -> adaptive pool ->
% Linear -> tanh; the dense layer holds 8.39 M of 8.6 M parameters, so it is
% close to a noise-conditioned MLP, not a waveform-synthesis architecture
% (WaveGAN/HiFi-GAN upsample). Zhou's citations: Mirza & Osindero (cGAN),
% Goodfellow (GAN), Fre-GAN (normalisation only), GAN-TTS (multi-task D),
% Saarinen & Koivunen (baseline). No other architecture needs citing.
Zhou trains the generator adversarially to imitate the victim's waveform; ours starts from random weights and learns only from the objective below, so the jammer resembles the victim only where resemblance pays.
Imitation as a starting point biases the generator toward the victim's modulation and makes it markedly easier to detect (Sec.~\ref{sec:exphist}).
Its modulation-label input is constant here, so the conditioning path is inactive.
\rar{Possible extension: condition on the target detector instead of the modulation, so one generator serves all detectors and shows whether detector-aware training generalizes. Depends on whether the MARL direction is pursued (README Sec. 4.3).}
% With one class the label embedding is a constant input (a learned bias).

\subsection{Training Against the Detector}
\label{sec:method:objective}
Starting from random weights, we minimize a differentiable form of Eq.~\eqref{eq:problem} by gradient descent, back-propagated through the link and the frozen detector; no discriminator is trained.
% The gradient is of Eq. loss (the surrogate), not of Eq. problem itself.
% run003 = --init random --steps 4000 (README Sec. 3.3f).
Two quantities in Eq.~\eqref{eq:problem} are replaced by differentiable surrogates.
First, $\mathbb{E}[\mathrm{BER}]$ is numerically zero wherever the jammer is weak, so we maximize $\log\mathbb{E}[\mathrm{BER}]$, computed exactly over the receiver noise: a bit whose noiseless matched-filter sample lies at signed distance $m_b$ from its decision boundary is in error with probability $Q(m_b/\sigma)$, with $\sigma^2=N_0/2$ the per-axis noise variance, and the logarithm of the mean over bits is evaluated with a log-sum-exp, so it never underflows.
% attacks.ber_logprob / log_expected_ber: log_ndtr(-m / sqrt(N0/2)), float64.
% E[BER] at SNR 30 dB underflows for a weak jammer (Q(~76) per bit) -- the
% wall that stalled the scalar-reward attempts (README A.5, 3.3e).
Second, with $\psi(r)$ the target detector's statistic, $\tau_\alpha$ its threshold, and $s_\psi$ the standard deviation of $\psi$ on clean frames, $P_{\mathrm{det}}$ is replaced by
\begin{equation}
\tilde P_{\mathrm{det}}=\mathbb{E}\Bigl[\bigl(1+e^{-(\psi(r)-\tau_\alpha)/s_\psi}\bigr)^{-1}\Bigr],
\label{eq:softpdet}
\end{equation}
which tends to the hard detection probability as $s_\psi\to0$.
% detectors.soft_pdet.
Eq.~\eqref{eq:softpdet} is the generator loss of a GAN whose only discriminator is the defender, separating clean from jammed received frames; here that discriminator is frozen, and indistinguishability is traded against damage instead of being pursued alone.
% User 2026-09-26: this replaces Zhou's imitation discriminator, which asks
% "does the jammer look like the victim's signal?" -- a question nobody in the
% threat model asks -- with the one the defender asks. Co-adaptive version:
% Future Work, parked below.
The generator $G$ minimizes
\begin{equation}
\mathcal{L}(G)=-\log\mathbb{E}[\mathrm{BER}]+\beta\,\tilde P_{\mathrm{det}},
\label{eq:loss}
\end{equation}
which is the objective of Eq.~\eqref{eq:problem} with effectiveness on a logarithmic scale, so $\beta$ takes a different numerical range here.
The power constraint of Eq.~\eqref{eq:power} is applied inside the forward pass: it removes the generator's overall gain, so the gradient acts only on the shape of the waveform, and power never enters the loss.
% channel.receive -> jammers.scale_to_jsr; verify.py 14b/14c (finite, nonzero
% gradient on all of G's parameters through Sionna's ApplyTimeChannel, the
% matched filter and the detector).
At each step the received JSR is drawn uniformly in decibels from a band that starts just below the JSR at which the target detector begins to flag, because where the detector flags every frame the logistic function in Eq.~\eqref{eq:softpdet} is flat and the detection term carries no gradient.
% The per-target band is load-bearing: on the power band the CNN's sigmoid'
% is ~1e-11 and no beta recovers a gradient (README Sec. 3.1 trap 1).
One generator is trained per target detector and value of $\beta$.

\subsection{Differentiating Through the Deployed Detector}
\label{sec:method:ste}
The energy and kurtosis statistics are differentiable as they stand.
The learned detector is not: its input image maps the spectrogram onto a 256-level colour map, a staircase whose gradient is zero.
We therefore use a straight-through estimator~\cite{bengio2013estimating}: the forward pass keeps the exact quantized image, so the deployed statistic is unchanged, and the backward pass substitutes the colour map's local slope.
The attacker thus back-propagates through the deployed network itself, whose weights stay frozen, with no surrogate network and no transfer gap.
% Weights frozen: detectors.load_cnn sets requires_grad False; the gradient
% flows through them to the input. detectors._viridis_apply(grad=True);
% verify.py 14d checks the forward values agree. Known and measured: the
% deployed CNN runs under fp16 autocast, D2 trains in fp32; decision-level
% equivalence asserted (README Sec. 3.1).
% The hard-decision reporting rule and the old Sec. method:design ("What the
% Comparisons Isolate") moved to Experiment Setup 2026-09-26 -- parked below.

% =============================================================================
%  PARKED FOR OTHER SECTIONS  (the methodology holds only the method)
% =============================================================================
%
% ---- Experiment Setup / Generator  (the architecture, moved out of
%      Methodology 2026-09-26 by the reframe; place before the recipe values)
% The generator concatenates $\mathbf{z}\in\mathbb{R}^{400}$ with a learned
% 16-dimensional embedding of the modulation label and maps the result through
% three one-dimensional convolutional blocks (convolution, batch
% normalization, leaky ReLU, dropout, max-pooling) and a dense layer with
% $\tanh$ output to the I and Q components of a segment of $M=1024$ samples,
% about $8.6\times10^{6}$ parameters in all, initialized at random.
%   [cgan/models.py. Zhou states z 400, label 16, 3 blocks, 1024 I + 1024 Q;
%   channels 64/128/256, kernel 5, dropout 0.3, pool-to-16 are ours (README
%   Q7). Since 2026-09-27 the method uses the architecture only: Zhou's
%   discriminator and imitation losses (adversarial + GP + classification;
%   feature matching, STFT distance, I/Q distribution distance; run001_G =
%   non-saturating BCE + GP, n_critic 1, all weights 1, MSE FM, log1p STFT,
%   headroom 1.0, 10k iterations, batch 128, Adam 2e-4) belong to the
%   Experiment History appendix, where the imitation start is reported.]
%
% ---- Conclusion & Future Work (co-adaptive defender = a GAN between attacker
%      and defender; user 2026-09-26) ---------------------------------------
% Retraining the defender on each round's jammer turns the procedure into a
% GAN between attacker and defender, with the defender as the discriminator;
% how the pair then co-adapts, and whether it converges, is future work.
%   [Supervisor axis "how defenders and attackers interact" (README Sec. B.2).
%   LITERATURE CHECK BEFORE ANY NOVELTY CLAIM: Shi, Davaslioglu & Sagduyu,
%   "Generative Adversarial Network for Wireless Signal Spoofing" (ACM WiseML
%   2019) already trains a generator against the receiver's classifier as
%   discriminator. Also README Sec. 4.1 #5 (co-adaptive defender, open).]
%
% ---- Conclusion & Future Work (team synchronization; promote to a
%      Methodology subsection + figure only if time allows, user 2026-09-26)
% Extended to a team, a first measurement without learning shows that $N_J$
% generative jammers match one strong jammer at lower per-jammer power only
% while their mutual timing error stays below about one symbol, and that
% uncompensated propagation delays already sit at that threshold; learning
% this synchronization is future work.
%   [README Sec. 3.3h, 3.3i (D4a): K = 4 aligned = one jammer at 7.3 dB less
%   per drone; knee sigma ~ 0.5-1 symbol, floor by 8-16; geometric offset
%   1.1 symbols rms on the 50 test drops. Transfer, not retrained. MUST state
%   the ceiling: a team cannot beat one jammer that controls its received
%   signal (containment) -- never "coordination beats a single jammer".]
%
% ---- Experiment Setup / Baselines: what each comparison isolates  (was
%      Methodology Sec. method:design, moved 2026-09-26, user; place after the
%      baseline definitions and the shaped control's recipe) ----------------
% The generator trained with $\beta=0$ shares everything with the
% detector-aware ones except the detection term, so their difference is the
% effect of detector awareness alone.
%   [2026-09-27: the plain (imitation) generator is no longer a main-text
%   comparator; imitation vs damage training is in the appendix.]
%   [CORRECTED 2026-09-26: the old sentence credited D1 -> D2 to "detector-aware
%   training alone", but D1 -> D2 also adds the BER term. At matched BER 3e-4,
%   30 dB, CNN P(det), run002 (4000 steps): D1 1.000 -> beta=0 0.579 -> CNN
%   beta=1 0.413; at 15 dB (E2) 0.94 -> 0.09 -> 0.09. D1 -> beta=0 is the big
%   effect (13.9 dB less power); the detector term adds -8..-20 pp vs the CNN
%   at 25-35 dB only; power/kurtosis targets are no better than beta=0 on
%   their own detector. README Sec. 3.3f "run002", 3.3g "E2 on run002".]
% The shaped-noise control is optimized for the same problem,
% Eq.~\eqref{eq:problem}, against the same detectors, but black-box over 48
% parameters rather than by gradient over the raw I/Q of a network of about
% $8.6\times10^{6}$ parameters, so its difference from the detector-aware
% generator reflects the signal family together with the optimizer each
% family admits.
% The fixed baselines bracket both, from the barrage floor to the
% unrealizable omniscient ceiling.
% Every reported number is measured with the deployed detector's hard
% decisions at its calibrated threshold, never with the surrogate of
% Eq.~\eqref{eq:softpdet}.
%   [Corrected 2026-09-26: the old text said "about 10^6 parameters" (it is
%   8,597,136) and "trained on the same objective". Same PROBLEM, different
%   surrogate: shaped = Eq. problem directly (linear E[BER], hard P_det,
%   beta {0, 0.01, 1}, one theta per JSR); GAN = Eq. loss (log E[BER], soft
%   P_det, beta {1..1e6}, one G per (target, beta) over a JSR band). The two
%   beta ranges are NOT comparable -- never put them on one axis.]
%
% ---- Experiment Setup / Baseline attackers: how the shaped control is trained
%      (goes right after the shaped jammer's definition, which is parked at the
%      end of sec_system_model_waveform.tex) -------------------------------
% The shaped jammer's parameter $\theta\in\mathbb{R}^{48}$ is optimized
% black-box by the covariance matrix adaptation evolution strategy
% (CMA-ES)~\cite{hansen2016cma}, started at $\theta=0$, the barrage jammer.
% It observes the target detector only through per-frame scores and decisions.
% Candidates are ranked by the objective of Eq.~\eqref{eq:problem} itself, with
% $\mathbb{E}[\mathrm{BER}]$ computed exactly as in Sec.~\ref{sec:method:objective}
% and $P_{\mathrm{det}}$ the hard fraction of flagged frames; ties, which occur
% wherever every candidate is flagged in every frame, are broken by the
% detector's mean score, weighted below the objective's resolution.
% One $\theta$ is optimized per target detector, value of $\beta$, and JSR, and
% a second, independent seed reaches the same operating points, so the control
% measures the family rather than the optimizer.
%   [train_shaped.py: pycma, sigma0 = 1 in log-power units, 150 generations,
%   256 frames/evaluation, common random numbers per generation; beta in
%   {0, 0.01, 1}; JSR -40:4:0 dB; tie weight 1e-6 on a monotone map of the
%   soft miss rate; seed check 6/6 pairs agree, |dP_det| <= 0.009 (README
%   Sec. 3.3e). Not RL: a static parameter vector is not a sequential decision.]
%
% ---- Experiment Setup (evaluation protocol) ---------------------------------
% Each attacker is swept over received JSR; at every point BER and each
% detector's P_det are measured on independent frames (fresh bits, jammer draw
% and noise per frame). A point selected from the sweep is re-measured on fresh
% frames before it is reported, because picking a maximum under a noisy P_det
% favours lucky draws.
%   [512 frames/point (binomial SE +-1.8 pp; differences < ~4 pp are not
%   meaningful), confirmation on 4096 fresh frames, pass = P_det <= alpha +
%   2 sigma (baselines.confirm). Matched-BER points interpolated on the 1 dB
%   JSR grid, linear in log BER; headline at BER 3e-4 (README Sec. 3.3f).]
% Noise ablation: SNR 0:5:40 dB plus a noiseless anchor; every threshold
% re-calibrated on clean frames at each level; the CNN keeps its 30 dB weights
% and only its colour scale is re-fitted; generators are evaluated, not
% retrained, so the ablation measures transfer.
%   [20 000 clean frames per level; README Sec. 3.3g setup + caveats; the
%   noiseless anchor is unusable for the CNN (OOD) and saturates both power
%   detectors -- an anchor, not a result.]
% Learned detector: EfficientNet-B0 as in~\cite{li2022jamming_ofdm}, their
% architecture retrained on this link with their four jammer types adapted to a
% single carrier, received JSR uniform in dB over [-20, +10]; 99.94 % accuracy
% at +10 dB.  [train_spectrogram_cnn.py; job 2266131, 100 epochs.]
